% 核心结论：删除说明移离斜线并写明有限电阻前提
% 草图：[真实几何与节点 | 条件、推导、结论]；17cm印宽，最小字号8pt。
% 图契约：等势支路的电流为零，才可以删去；图中的端口、几何、箭头和方程对应本题。
% 版面：沿用前三模块的几何图+推导分区；目标印宽17cm，简图至中等密度。
% 布局：左侧真实结构，右侧关键关系；真实电路边不作为架构流程箭头处理。
\documentclass[tikz,border=7pt]{standalone}
\usepackage{fontspec,xeCJK,amsmath,amssymb}
\IfFontExistsTF{Times New Roman}{\setmainfont{Times New Roman}}{\setmainfont{TeX Gyre Termes}}
\IfFontExistsTF{Heiti SC}{\setCJKmainfont[AutoFakeBold=1.5]{Heiti SC}}{\setCJKmainfont[AutoFakeBold=1.5]{FandolHei-Regular}}
\usetikzlibrary{arrows.meta,bending,calc,patterns,positioning,decorations.pathmorphing}
\definecolor{ink}{HTML}{243745}
\definecolor{blue}{HTML}{275DAB}
\definecolor{green}{HTML}{177C63}
\definecolor{red}{HTML}{BD4935}
\definecolor{orange}{HTML}{AA6C16}
\definecolor{purple}{HTML}{715393}
\tikzset{every node/.style={font=\fontsize{9.3}{12.5}\selectfont,text=ink},
arrow/.style={-{Stealth[length=5pt 1.5,width'=0pt 3]},shorten >=2pt,shorten <=1pt,line width=1pt},
arrow short/.style={-{Stealth[length=3pt 1.0,width'=0pt 2.5]},shorten >=1pt,shorten <=0pt,line width=.8pt},
axis/.style={arrow,draw=ink!60,line width=.6pt},
vec/.style={arrow,draw=blue},
wire/.style={draw=ink!75,line width=.8pt},
guide/.style={dashed,draw=ink!45,line width=.5pt},
card/.style={rounded corners=2pt,draw=blue!65,fill=blue!4,align=center,text width=6.5cm,inner xsep=9pt,inner ysep=7pt},
title/.style={font=\fontsize{10.5}{14}\selectfont\bfseries,text width=15cm,align=center},
label/.style={font=\fontsize{8.4}{11}\selectfont,align=center},
note/.style={label,text width=6.1cm}}
\begin{document}
\begin{tikzpicture}[x=1cm,y=1cm]
\path[use as bounding box] (-8,-3.5) rectangle(8,3.65);
\node[title] at (0,3.25) {\quad 等势支路的电流为零，才可以删去};
\begin{scope}[shift={(-4.05,0)}]
\draw[wire,draw=ink](-2,1.2)--(-1.2358105545709497,0.7414863327425698) (-0.7641894454290503,0.4585136672574302)--(0,0);\begin{scope}[shift={(-1.0,0.6)},rotate=-30.96375653207352]\draw[draw=ink,fill=white,line width=.8pt](-0.275,-.085)rectangle(0.275,.085);\end{scope}
\node[label]at(-0.7839117827204389,0.9601470287992685){$R_1$};
\draw[wire,draw=ink](0,0)--(0.7641894454290503,0.4585136672574302) (1.2358105545709497,0.7414863327425698)--(2,1.2);\begin{scope}[shift={(1.0,0.6)},rotate=30.96375653207352]\draw[draw=ink,fill=white,line width=.8pt](-0.275,-.085)rectangle(0.275,.085);\end{scope}
\node[label]at(0.7839117827204389,0.9601470287992685){$R_2$};
\draw[wire,draw=ink](-2,-1.2)--(-1.2358105545709497,-0.7414863327425698) (-0.7641894454290503,-0.4585136672574302)--(0,0);\begin{scope}[shift={(-1.0,-0.6)},rotate=30.96375653207352]\draw[draw=ink,fill=white,line width=.8pt](-0.275,-.085)rectangle(0.275,.085);\end{scope}
\node[label]at(-1.216088217279561,-0.23985297120073146){$R_3$};
\draw[wire,draw=ink](0,0)--(0.7641894454290503,-0.4585136672574302) (1.2358105545709497,-0.7414863327425698)--(2,-1.2);\begin{scope}[shift={(1.0,-0.6)},rotate=-30.96375653207352]\draw[draw=ink,fill=white,line width=.8pt](-0.275,-.085)rectangle(0.275,.085);\end{scope}
\node[label]at(1.216088217279561,-0.23985297120073146){$R_4$};
\fill[ink] (0,0)circle(2pt);
\draw[guide,red](2,1.2)--(2,-1.2);\node[label,fill=white]at(2.65,0){$V_i=V_j$};\fill[ink] (2,1.2)circle(2pt);\node[above=4pt] at(2,1.2){$i$};
\fill[ink] (2,-1.2)circle(2pt);\node[below=4pt] at(2,-1.2){$j$};
\node[note]at(0,-2.8){虚线表示待删电阻支路，不表示理想导线};
\end{scope}
\node[card,text width=6.5cm,draw=blue!65,fill=blue!4]at(3.9,1.8){$V_i=V_j\quad\Rightarrow\quad I_{ij}=0$};
\node[card,text width=6.5cm,draw=green!65,fill=green!4]at(3.9,0){删去该有限电阻支路不改变本次实验};
\node[card,text width=6.5cm,draw=orange!65,fill=orange!4]at(3.9,-1.8){先用对称性或方程证明等势\\端口改变后需要重新判断};

\end{tikzpicture}\end{document}